\section{Conclusion}
The response information needed to procure a dynamic service is determined by the institution, not by the accuracy of an arbitrary policy prediction. Even complete coordinate-value information can leave the procurement decision unidentified. A purchaser-directed joint query can resolve that ambiguity, with an exact noise threshold and an explicit value of information. This result explains why joint response sets, rather than separate coordinate extrema, are the appropriate output of the value-verification layer.

Implementation has a similarly precise behavioral content. Uniform state-wise enforcement and initial-state enforcement are distinct objects. Reachability or a comparison of adjustment gains can connect them; an explicit initial certificate can establish the ordering without those primitive conditions. Fee slack then prices globally near-optimal surrender through a sharp loss-budget bound. Whole-interval action certificates remain valid at positive tolerance only when their losses are combined through occupation probabilities rather than screened by unrelated local cutoffs.

The applications provide complementary evidence. The Brownian benchmark completes a continuous-model verification and participation-adjusted procurement comparison with analytic candidates and rational numerical constants. The original settlement economy retains its operating primitives, initial financial mandates, full charge continuum, institution counterfactuals, and adverse mechanism cases. Its newly bracketed initial thresholds establish a separate adjustment effect on the same frozen arrays. The continuous benchmark does not disguise the remaining constructor-inclusion question for those arrays.

Neural Bellman Operators retain their role as an approximation and proposal architecture connected to these economic messages. Their useful predictions are evaluated, their errors are charged, and their unfavorable outcomes remain in the resource comparison alongside evaluated teacher banks and exact dynamic programming. The distinction between proposing a policy, verifying a value, identifying a response, and choosing an executable institution is what permits the economic conclusions to survive both ties and imperfect approximation.
